\section{Equivariant kernels and their multiplicity spaces}
\label{l:equivariant-section}

Partial-permutation information can be imposed on a transition kernel
without first selecting a starting permutation.  This preserves a group
action on the entire state space.  The resulting Fourier decomposition
has two parts: an irreducible carrier, on which the kernel acts as the
identity, and a multiplicity space, on which it may act nontrivially.
The latter space can be large even when the carrier is one-dimensional.
We prove two amplification results that retain these multiplicities.

Throughout this section kernels act on functions by
\[
 (Af)(x)=\sum_y A(x,y)f(y).
\]
All Hilbert--Schmidt norms use counting measure and the ordinary trace.
For a finite state space $\Omega$, let $U_\Omega$ be the matrix with
every entry $1/|\Omega|$.  Thus $U_\Omega$ is the orthogonal projection
onto the normalized constant vector.  A nonnegative matrix with row and
column sums at most one is a contraction on $\ell^2(\Omega)$: indeed,
\[
 \sum_x\left|\sum_y A(x,y)f(y)\right|^2
 \le \sum_{x,y}A(x,y)|f(y)|^2\le\sum_y|f(y)|^2.
\]
We will use this fact for kernels obtained by deleting entries of a
doubly stochastic kernel.

\subsection{Sixteen separately truncated factors}

Let $A$ and $V$ be two sets of size $n$, and let
$\Omega=\operatorname{Bij}(A,V)$ be the set of bijections from labels
to positions.  Partition $A$ into sixteen blocks $A_1,\ldots,A_{16}$,
of sizes $m_i$, and write
\[
 H_i=S_{A_i},\qquad H=\prod_{i=1}^{16}H_i.
\]
These groups act on $\Omega$ on the right by composition:
$x\cdot h=x\circ h$.  The quotient $X_i=\Omega/H_i$ records the
positions of every label outside $A_i$.  Its size is
$L_i=n!/m_i!$, and each of its fibers has $m_i!$ elements.

If a kernel $K$ obeys $K(xh,yh)=K(x,y)$ for $h\in H$, its quotient
kernel on $X_i$ is well defined by
\[
 q_i(a,b)=\sum_{y:\,\pi_i(y)=b}K(x,y),\qquad \pi_i(x)=a,
\]
where $\pi_i:\Omega\to X_i$ is the quotient map.  Equivariance under
$H_i$ makes the right side independent of the representative $x$.
Fix a number $\varepsilon>0$ and define a retained kernel by
\begin{equation}
 T_i(x,y)=K(x,y)\,
 \mathbf1\left\{q_i(\pi_i(x),\pi_i(y))
                         \le\frac{1+\varepsilon}{L_i}\right\}.
 \label{l:sixteen-cutoff}
\end{equation}
Since $H_i$ is normal in $H$, the quotient has a right $H$-action.
Equivariance of $K$ gives $q_i(ah,bh)=q_i(a,b)$ for $h\in H$.
The cutoff therefore preserves the action of all of $H$.

We specify one further hypothesis on $T_i$.  In the orthogonal
$H_i$-isotypic decomposition
\[
 \ell^2(\Omega)=\bigoplus_{\alpha\in\widehat H_i}
                 V_\alpha\otimes\mathcal M_{i,\alpha},
 \qquad
 T_i=\bigoplus_\alpha I_{V_\alpha}\otimes T_i^\alpha,
\]
$\mathcal M_{i,\alpha}$ is the full multiplicity space, including all
quotient fibers.  In particular $T_i^{\mathrm{sgn}}$ denotes the
operator on the sign multiplicity space, not a single scalar Fourier
coefficient.

\begin{theorem}[Sixteen equivariant subkernels]
\label{l:equivariant-sixteen}
Consider a sequence of the preceding state spaces and partitions with
$\min_i m_i\to\infty$.  Let $K$ be a doubly stochastic,
$H$-equivariant kernel.  Suppose that, for every $i$,
\begin{equation}
 \frac1{L_i}\sum_{a\in X_i}
 \frac12\sum_{b\in X_i}|q_i(a,b)-L_i^{-1}|\le\delta,
 \label{l:sixteen-marginal-input}
\end{equation}
where $\delta\to0$.  Choose $\varepsilon\to0$ with
$\delta/\varepsilon\to0$, form the kernels
\eqref{l:sixteen-cutoff}, and assume also that
\begin{equation}
 \beta_i:=\|T_i^{\mathrm{sgn}}\|_{\mathrm{HS}}^2\longrightarrow0
 \qquad(1\le i\le16).
 \label{l:sixteen-sign-input}
\end{equation}
Then
\[
 \frac1{|\Omega|}\sum_{x\in\Omega}
       \|K^{16}(x,\cdot)-U_\Omega(x,\cdot)\|_{\mathrm{TV}}
 \longrightarrow0.
\]
If $K$ is equivariant under a transitive group of permutations of
$\Omega$, the same conclusion holds for the maximum over starting
states.

More precisely, put $T=T_1\cdots T_{16}$.  For a product type
$\boldsymbol\alpha=(\alpha_1,\ldots,\alpha_{16})$, set
$f_i=\dim\alpha_i$ and $D=\prod_i f_i$.  If
$T_{i,\boldsymbol\alpha}$ is the operator on its full
$H$-multiplicity space, then
\begin{equation}
 \|T_{i,\boldsymbol\alpha}\|_{\mathrm{HS}}^2
       \le(1+\varepsilon)\frac{f_i^2}{D},\qquad
 D\|T_{1,\boldsymbol\alpha}\cdots
             T_{16,\boldsymbol\alpha}\|_{\mathrm{HS}}^2
       \le(1+\varepsilon)^{16}D^{-13}.
 \label{l:sixteen-type-bounds}
\end{equation}
\end{theorem}

\begin{proof}
We first bound the deleted mass, then use each of the sixteen quotient
caps on its corresponding factor.  Write
\[
 \gamma_i=\frac1{|\Omega|}\sum_{x,y}(K(x,y)-T_i(x,y)).
\]
On an atom with $q_i(a,b)>(1+\varepsilon)/L_i$, its positive excess
over uniform is at least
$\varepsilon q_i(a,b)/(1+\varepsilon)$.  Summing over such atoms and
then averaging over $a$ gives
\begin{equation}
 \gamma_i\le\frac{1+\varepsilon}{\varepsilon}\delta=o(1).
 \label{l:sixteen-loss}
\end{equation}
Every $T_i$ is row- and column-substochastic.  A uniform starting
distribution remains dominated by uniform after any product of these
kernels.  The average mass lost in $T$ is consequently at most
$\Gamma=\sum_i\gamma_i$.  Also $0\le T\le K^{16}$ entrywise.
The triangle inequality followed by Cauchy--Schwarz over all
$|\Omega|^2$ entries gives
\begin{equation}
 \frac1{|\Omega|}\sum_x
       \|K^{16}(x,\cdot)-U_\Omega(x,\cdot)\|_{\mathrm{TV}}
 \le \frac\Gamma2+\frac12\|T-U_\Omega\|_{\mathrm{HS}}.
 \label{l:sixteen-average-tv}
\end{equation}

Fix $i$ and first decompose only under $H_i$.  Each quotient fiber
is a regular $H_i$-set.  After choosing one representative per fiber,
the block of $T_i$ between fibers $a,b$ is convolution by nonnegative
weights of total mass
\[
 t_i(a,b)=q_i(a,b)
       \mathbf1\{q_i(a,b)\le(1+\varepsilon)/L_i\}.
\]
On the multiplicity space of a regular-fiber type $\alpha$, of degree
$f_\alpha$, this block is a weighted sum of unitary $f_\alpha$ by
$f_\alpha$ matrices.  Its Hilbert--Schmidt norm is at most
$\sqrt{f_\alpha}\,t_i(a,b)$.  Since each row of $t_i$ has mass at
most one, the cap gives
\begin{equation}
 \|T_i^\alpha\|_{\mathrm{HS}}^2
 \le f_\alpha\sum_{a,b}t_i(a,b)^2
 \le(1+\varepsilon)f_\alpha.
 \label{l:sixteen-single-hs}
\end{equation}

Now decompose under the whole product $H$.  Its
$\boldsymbol\alpha$-isotype is
\[
 \left(\bigotimes_{j=1}^{16}V_{\alpha_j}\right)
                  \otimes\mathcal M_{\boldsymbol\alpha}.
\]
Viewed only as an $H_i$-space, it contributes $D/f_i$ copies of
$\mathcal M_{\boldsymbol\alpha}$ to the multiplicity of $\alpha_i$.
Thus \eqref{l:sixteen-single-hs} implies
\[
 \frac D{f_i}\|T_{i,\boldsymbol\alpha}\|_{\mathrm{HS}}^2
 \le(1+\varepsilon)f_i.
\]
Multiplying these inequalities and using Hilbert--Schmidt
submultiplicativity proves \eqref{l:sixteen-type-bounds}.

For $u>0$, write
$\mathcal Z_u(m)=\sum_{\alpha\in\widehat S_m}(\dim\alpha)^{-u}$.
The degree estimate of Theorem~\ref{l:degree-all-powers} gives
$\mathcal Z_{13}(m)\to2$.  The total squared Hilbert--Schmidt contribution
of all product types with $D>1$ is therefore at most
\begin{equation}
 (1+\varepsilon)^{16}
       \left(\prod_{i=1}^{16}\mathcal Z_{13}(m_i)-2^{16}\right)=o(1).
 \label{l:sixteen-nonscalar-sum}
\end{equation}
The subtracted types are precisely the products of trivial and sign
representations, since eventually every $m_i\ge2$.

It remains to handle these $2^{16}$ scalar carrier types; their
multiplicity spaces have not been discarded.  If a product type has
a sign coordinate $i$, its $T_i$ block has squared Hilbert--Schmidt
norm at most $\beta_i$ by \eqref{l:sixteen-sign-input}.  All the
other factors are operator contractions.  Hence its product block
has squared Hilbert--Schmidt norm at most $\beta_i$.  The sum over
these finitely many types tends to zero.

For the all-trivial type, let $T_{i,0}$ be its multiplicity operator
and let $u=|\Omega|^{-1/2}\mathbf1$ be the constant vector in that
space.  The restriction of $U_\Omega$ is $P_u=uu^*$.
Equation~\eqref{l:sixteen-single-hs} for the trivial $H_i$ type gives
$\|T_{i,0}\|_{\mathrm{HS}}^2\le1+\varepsilon$, whereas
\[
 \langle u,T_{i,0}u\rangle=1-\gamma_i.
\]
Consequently
\begin{equation}
 \|T_{i,0}-P_u\|_{\mathrm{HS}}^2
 =\|T_{i,0}\|_{\mathrm{HS}}^2+1
       -2\langle u,T_{i,0}u\rangle
 \le\varepsilon+2\gamma_i.
 \label{l:sixteen-trivial-rank-one}
\end{equation}
Telescoping the product against $P_u^{16}=P_u$, with operator
contractions on either side of each difference, gives
\[
 \|T_{1,0}\cdots T_{16,0}-P_u\|_{\mathrm{HS}}
 \le\sum_{i=1}^{16}\sqrt{\varepsilon+2\gamma_i}=o(1).
\]
This is a bound on the entire all-trivial multiplicity space.

The three estimates prove $\|T-U_\Omega\|_{\mathrm{HS}}\to0$,
so \eqref{l:sixteen-average-tv} proves average mixing.  Finally,
equivariance under a transitive group makes all row distances of
$K^{16}$ from uniform equal.  That additional symmetry gives the
worst-start conclusion.
\end{proof}

The sign hypothesis concerns the retained factors themselves.  A
probabilistic construction can verify it by cancellation inside each
retained fiber block.  The quotient total-variation estimate alone
supplies the mass loss and the general type bound, but does not
control the sign multiplicity operator.

\subsection{Three groups and two-sided quotient control}

The next result uses only three label groups, at the cost of controlling
every quotient row and column.  All three cutoffs are imposed on one
kernel.  Fourier analysis then gives one matrix for each product type,
and fifteen powers of that matrix supply the summable dimension saving.

Let $\mathcal S$ be a nonempty finite set of size $M$, let
$H_i=S_{m_i}$ for $1\le i\le3$, and put
$H=H_1\times H_2\times H_3$ and $\Omega=\mathcal S\times H$.
The right $H$-action is $(s,g)h=(s,gh)$.  Its equivariant kernels
have the form
\begin{equation}
 K((s,g),(s',g'))=p_{s,s'}(g'g^{-1}).
 \label{l:three-convolution-kernel}
\end{equation}
Indeed simultaneous right multiplication reduces the source group
coordinate to the identity.  We always take the functions
$p_{s,s'}:H\to[0,\infty)$ to satisfy
\begin{equation}
 \sum_{s',h}p_{s,s'}(h)=1\quad(s\in\mathcal S),\qquad
 \sum_{s,h}p_{s,s'}(h)=1\quad(s'\in\mathcal S).
 \label{l:three-bistochastic}
\end{equation}
These are exactly the row and column normalizations of $K$.
For $H_{-i}=\prod_{j\ne i}H_j$, define
\[
 q_i(s,s',h_{-i})=\sum_{h_i\in H_i}p_{s,s'}(h),
 \qquad M_i=M|H_{-i}|.
\]
The associated quotient state consists of $s$ and the two retained
group coordinates.  Its uniform transition probability is $1/M_i$.
Transposition of $K$ replaces the fibers by
$\bar p_{s',s}(h)=p_{s,s'}(h^{-1})$.  Thus a quotient row estimate
for the transpose is exactly the quotient column estimate below,
after inversion of $h_{-i}$.

\begin{theorem}[Three-group two-sided truncation]
\label{l:three-group}
Consider a sequence of kernels \eqref{l:three-convolution-kernel}
satisfying \eqref{l:three-bistochastic}, with
$\min_i m_i\to\infty$.  Suppose $\delta\to0$ and, for every $i$,
\begin{align}
 \sum_{s',h_{-i}}|q_i(s,s',h_{-i})-M_i^{-1}|&\le2\delta
                                      &&(s\in\mathcal S),\label{l:three-row-input}\\
 \sum_{s,h_{-i}}|q_i(s,s',h_{-i})-M_i^{-1}|&\le2\delta
                                      &&(s'\in\mathcal S).\label{l:three-column-input}
\end{align}
Let $\chi(h)=\prod_{i=1}^3\mathrm{sgn}(h_i)$, and assume that the
untrimmed all-sign matrix
\[
 A_\chi(s,s')=\sum_{h\in H}p_{s,s'}(h)\chi(h^{-1})
\]
has absolute row and column sums at most $a$, where $a\to0$.
Then $K^{15}$ mixes in average row total variation:
\[
 \frac1{|\Omega|}\sum_x
       \|K^{15}(x,\cdot)-U_\Omega(x,\cdot)\|_{\mathrm{TV}}
 \longrightarrow0.
\]
The maximum row distance also tends to zero if $K$ is equivariant
under a transitive group of permutations of $\Omega$.

Specifically, retain the fibers
\begin{equation}
 \widetilde p_{s,s'}(h)=p_{s,s'}(h)
     \prod_{i=1}^3\mathbf1\{q_i(s,s',h_{-i})\le2/M_i\},
 \label{l:three-cutoff}
\end{equation}
and let $\widetilde K$ be the corresponding kernel.  At most
$6\delta$ mass is removed from every row and column.  For an
irreducible product type $\rho=\rho_1\otimes\rho_2\otimes\rho_3$,
put $D_i=\dim\rho_i$ and $D=\prod_iD_i$.  Its multiplicity matrix
is the $M$ by $M$ block matrix
\begin{equation}
 B_\rho(s,s')=\sum_{h\in H}\widetilde p_{s,s'}(h)\rho(h^{-1}),
 \label{l:three-fourier-block}
\end{equation}
with blocks of size $D$ by $D$, and it satisfies
\begin{equation}
 \|B_\rho\|_{\mathrm{HS}}^2
       \le4\frac{D_i}{\prod_{j\ne i}D_j}\quad(1\le i\le3),
 \qquad
 D\|B_\rho^{15}\|_{\mathrm{HS}}^2\le2^{30}D^{-4}.
 \label{l:three-type-bounds}
\end{equation}
\end{theorem}

\begin{proof}
For fixed $s$, the entries of each $q_i$ sum to one.  The same is
true for fixed $s'$ by \eqref{l:three-bistochastic}.  Therefore the
positive excess over uniform in either a quotient row or a quotient
column is at most $\delta$.  On an atom exceeding twice uniform,
its mass is at most twice its positive excess.  Each cutoff loses
at most $2\delta$ in every row and column.  Taking the union of the
three deleted sets proves the loss bound $6\delta$.
The retained kernel is row- and column-substochastic, so its blocks
are operator contractions.  Also its $i$th quotient marginal is
everywhere bounded by $2/M_i$: an original atom above that value
has been entirely deleted.

We check the Fourier orientation before estimating the blocks.  The
unitary regular decomposition of the right $H$-action has carrier
$V_\rho$ and multiplicity $D$ on each fiber $s$.  In the
matrix-coefficient basis $\rho(g^{-1})_{ab}$, left multiplication
$g\mapsto hg$ replaces
\[
 \rho(g^{-1})\quad\hbox{by}\quad
 \rho(g^{-1}h^{-1})=\rho(g^{-1})\rho(h^{-1}).
\]
For each fixed carrier index $a$, the index $b$, together with $s$,
is therefore acted on by \eqref{l:three-fourier-block}.  Each such
matrix occurs $D$ times.  In particular composition of the deck
kernels gives powers of $B_\rho$, with no adjoint or reversal of
the factors.

Fix $s,s'$ and an index $i$.  First sum over its group coordinate:
\[
 F(h_{-i})=\sum_{h_i}\widetilde p_{s,s'}(h)\rho_i(h_i^{-1}).
\]
The quotient cap and unitarity give
$\|F(h_{-i})\|_{\mathrm{HS}}\le2\sqrt{D_i}/M_i$.
Applying Plancherel entry by entry on $H_{-i}$ yields
\begin{align*}
 &\sum_{\sigma\in\widehat H_{-i}}(\dim\sigma)
 \left\|\sum_{h_{-i}}\sigma(h_{-i}^{-1})\otimes
                       F(h_{-i})\right\|_{\mathrm{HS}}^2\\
 &\hspace{25mm}=|H_{-i}|\sum_{h_{-i}}\|F(h_{-i})\|_{\mathrm{HS}}^2
 \le\frac{4D_i}{M^2}.
\end{align*}
Keep the term $\sigma=\bigotimes_{j\ne i}\rho_j$ and sum over the
$M^2$ pairs $(s,s')$.  A unitary permutation of tensor factors
identifies the resulting matrix with $B_\rho$, proving the first
bound in \eqref{l:three-type-bounds}.  Choose an index with
$D_i\le D^{1/3}$.  Then
\[
 \|B_\rho\|_{\mathrm{HS}}^2\le4D^{-1/3},\qquad
 D\|B_\rho^{15}\|_{\mathrm{HS}}^2
 \le D(4D^{-1/3})^{15}=2^{30}D^{-4}.
\]
Hilbert--Schmidt submultiplicativity suffices here; no normality is
assumed.  By Theorem~\ref{l:degree-all-powers}, the sum over $D>1$
is at most
\[
 2^{30}\left(\prod_{i=1}^3 \mathcal Z_4(m_i)-8\right)=o(1).
\]

We have controlled every carrier of dimension greater than one.
The eight scalar carriers remain: one is all-trivial, six contain
both a trivial and a sign factor, and one is all-sign.  In every
case the corresponding $M$ by $M$ matrix has Hilbert--Schmidt norm
at most two, by \eqref{l:three-type-bounds}.

Consider first one of the six mixed types.  Choose an index $i$
where $\rho_i$ is trivial, so that its matrix entries are the signed
sums of the retained marginal $\widetilde q_i$.  Replacing that
marginal by $1/M_i$ gives zero, since at least one of the other two
factors is sign and its group sum vanishes.  In each row and column,
\[
 \sum|\widetilde q_i-M_i^{-1}|
 \le\sum|q_i-M_i^{-1}|+\sum(q_i-\widetilde q_i)
 \le8\delta.
\]
Thus $B_\rho$ has absolute row and column sums at most $8\delta$,
and the same Cauchy--Schwarz argument used for nonnegative kernels,
applied to absolute entries, gives
$\|B_\rho\|_{\mathrm{op}}\le8\delta$.  Hence
\[
 \|B_\rho^{15}\|_{\mathrm{HS}}
       \le2(8\delta)^{14}=o(1).
\]

For the all-trivial type, uniform replacement instead gives the
matrix $J$ with every entry $1/M$.  Write $B_0=J+E$.
The preceding row and column estimate gives
$\|E\|_{\mathrm{op}}\le8\delta$, while
$\|E\|_{\mathrm{HS}}\le\|B_0\|_{\mathrm{HS}}+\|J\|_{\mathrm{HS}}
\le3$.  Expand $(J+E)^{15}$ into ordered words.  The word containing
only $J$ equals $J$, because $J$ is an orthogonal projection.  Every
word containing both letters has a factor $J$ of Hilbert--Schmidt
norm one and at least one factor $E$ of operator norm at most
$8\delta$, so its Hilbert--Schmidt norm tends to zero.  The word
$E^{15}$ has norm at most $3(8\delta)^{14}$.  There are only $2^{15}$
words.  This proves
\[
 \|B_0^{15}-J\|_{\mathrm{HS}}=o(1)
\]
without bounding the size $M$ of the trivial multiplicity space.

Finally consider the all-sign type.  Its untrimmed matrix is
$A_\chi$.  Deleting at most $6\delta$ mass from each row and column
changes each corresponding absolute sum by at most $6\delta$.
Consequently
\[
 \|B_\chi\|_{\mathrm{op}}\le a+6\delta,
 \qquad
 \|B_\chi^{15}\|_{\mathrm{HS}}
       \le2(a+6\delta)^{14}=o(1).
\]
This is precisely where the separate all-sign hypothesis is used.

The uniform kernel has only the all-trivial block $J$.  Summing all
blocks with their carrier multiplicities $D$ now gives
$\|\widetilde K^{15}-U_\Omega\|_{\mathrm{HS}}=o(1)$.
For all sufficiently large members of the sequence $6\delta<1$.
Every row of $\widetilde K^{15}$ has mass at least
$(1-6\delta)^{15}$ and is dominated by the corresponding row of
$K^{15}$.  Thus, as in \eqref{l:sixteen-average-tv},
\[
 \frac1{|\Omega|}\sum_x
       \|K^{15}(x,\cdot)-U_\Omega(x,\cdot)\|_{\mathrm{TV}}
 \le45\delta+\frac12\|\widetilde K^{15}-U_\Omega\|_{\mathrm{HS}}
 =o(1).
\]
The final transitive-equivariance assertion follows because the
untrimmed kernel then has the same row distance at every state.
\end{proof}

The two-sided assumptions have distinct roles.  Quotient rows bound
the forward mass removed; quotient columns bound the reverse mass
removed and the operator norms of the scalar error matrices.  The
all-sign estimate is additional data.  In applications where it is
proved by cancellation at the end of a time block, the distribution
just before that canceling step must also be controlled in both
orientations.
